% TOP_PROBLEM: {"id": 1421, "title": "Oberwolfach Problem", "category_id": 3, "category": {"id": 3, "name": "graph_theory", "display_name": "Graph Theory", "description": "Problems involving graphs, networks, and their properties.", "slug": "graph-theory", "order_index": 3, "created_at": "2026-07-31T15:26:25.671Z"}, "status": "open"}
% FIRST_POSTED: null
% ENTRY_KIND: statement_only
% Imported from: batch07.tex; UnsolvedMath ID 1421
\documentclass[11pt]{article}
\usepackage[margin=25mm]{geometry}
\usepackage{fontspec}
\setmainfont{Latin Modern Roman}
\usepackage{amsmath,amssymb,mathtools}
\usepackage{unicode-math}
\setmathfont{Latin Modern Math}
\usepackage{xurl,hyperref}
\hypersetup{colorlinks=true,urlcolor=blue}
\setcounter{secnumdepth}{0}
\setlength{\parindent}{0pt}
\setlength{\parskip}{5pt}
\setlength{\emergencystretch}{3em}
\newcommand{\sref}[2]{\hyperlink{source:#1:#2}{[S#2]}}
\newcommand{\cataloguescope}[1]{\paragraph{Recorded target.}#1}
\begin{document}
% UnsolvedMath ID 1421; source catalogue code GRAPH-034
\setcounter{section}{1420}
\section{Oberwolfach Problem}\label{1421}
{\small\sffamily UnsolvedMath ID 1421 · Graph Theory}
\subsection{Definitions and mathematical statement}

\paragraph{Shared notation.}$\mathbb N=\{1,2,\ldots\}$, and $\mathbb Z,\mathbb Q,\mathbb R,\mathbb C$ denote the integers, rationals, reals and complex numbers. Any other conventions are specified locally.


Let $H$ be a finite simple undirected $2$-regular graph on $n$ vertices, meaning every vertex has degree two; thus $H$ is a disjoint union of cycles of lengths at least three. The complete graph $K_n$ has all unordered pairs of different vertices as edges. An $H$-decomposition of $K_n$ is a collection of spanning subgraphs $H_1,\ldots,H_r$, each isomorphic to $H$, such that their edge sets partition $E(K_n)$. Degree counting forces $2r=n-1$, so $n$ must be odd and $r=(n-1)/2$.
\paragraph{Statement recorded in the supplied source.}
Classify all finite $2$-regular graphs $H$ for which $K_{|V(H)|}$ admits an $H$-decomposition. Equivalently, for every odd $n\ge3$ and every partition $n=\ell_1+\cdots+\ell_s$ with $\ell_i\ge3$, determine whether $K_n$ can be decomposed into $(n-1)/2$ spanning copies of $C_{\ell_1}\sqcup\cdots\sqcup C_{\ell_s}$. Here $C_\ell$ is a cycle of length $\ell$, and $\sqcup$ denotes disjoint union. Even orders are impossible by parity; the finite classification includes the small exceptional cases, not only sufficiently large orders. \sref{1421}{2}
\subsection{Short English statement}
Determine precisely which collections of cycle lengths can form repeated spanning factors that partition all edges of a complete graph.
\subsection{Sources}
\begin{description}
\item[\hypertarget{source:1421:1}{S1}] ULAM.AI and the UnsolvedMath Contributors, UnsolvedMath: A Curated Collection of Open Mathematics Problems, record 1421 (GRAPH-034). CC BY 4.0. \url{https://huggingface.co/datasets/ulamai/UnsolvedMath}.
\item[\hypertarget{source:1421:2}{S2}] Peter Keevash and Katherine Staden, The Generalised Oberwolfach Problem, arXiv:2004.09937v1 (2020), Section 1, “Oberwolfach Problem (Ringel)”, page 1. \url{https://arxiv.org/pdf/2004.09937v1}.
\end{description}
\label{end:1421}
\end{document}
